\appendix
\section{Integration test catalogue}
\label{sec:IntegrationTestAppendix}

This appendix documents every test suite in the project that exercises
behaviour against a live Kafka broker (Docker-backed integration tests)
plus the fault detection and alarm-state unit test suites for completeness.
Docker-backed integration tests extend \texttt{KafkaIntegrationTestBase}, which
uses Testcontainers to spin up a real Kafka broker (Confluent CP 7.8.0 in KRaft
mode) per test class. Each integration test class creates its own topics with
randomised suffixes to avoid collisions and disposes of all resources in
\texttt{@After} / \texttt{@AfterClass} hooks.

The integration tests are run with:

\begin{verbatim}
mvn test -Dtest='*IntegrationTest'
\end{verbatim}

Timeouts default to 60~seconds and can be overridden with
\texttt{-Ddurga.it.timeout.seconds=N} or the environment variable
\texttt{DURGA\_IT\_TIMEOUT\_SECONDS}.

% ===========================================================================

\subsection{ProcessEventKafkaIntegrationTest}
\label{it:processEventKafka}

\textbf{Purpose:} Verify that \texttt{ProcessEvent} records survive a
Kafka produce-consume round-trip with full wire-format fidelity.

\textbf{Topic:} A single non-compacted topic created per test class.
The test uses a plain \texttt{KafkaProducer} and \texttt{KafkaConsumer} with
manual offset management (\texttt{enable.auto.commit=false}, \\
\texttt{auto.offset.reset=earliest}).

\textbf{Scenarios (3):}

\begin{description}

\item[shouldProduceAndConsumeSingleEvent] Publishes one \texttt{ProcessEvent}
with an \texttt{ErrorInfo} payload (\texttt{code="TASK\_FAILED"}). Consumes it
and asserts that the \texttt{processInstanceId}, \texttt{processId},
\texttt{eventType} and \texttt{error().code()} are preserved exactly.
Exercises the full \texttt{toJson} / \texttt{fromJson} pipeline over the wire.

\item[shouldProduceAndConsumeMultipleEvents] Publishes five distinct events
with monotonically increasing \texttt{processInstanceId}s, consumes them all
in an at-least-once poll, collects IDs, and asserts all five were received
and each has a valid \texttt{processId} and \texttt{eventType}.

\item[shouldRoundTripFromJsonToJson] Publishes a compact-constructor event
(no explicit \texttt{eventType}, relying on the \texttt{Status} inference),
consumes it, re-serializes the parsed result, parses again, and asserts that
the double-parsed event matches the original \texttt{processInstanceId},
inferred \texttt{eventType}, and nested \texttt{payload} values.

\end{description}

\textbf{Why it matters:} The entire monitoring model depends on
\texttt{ProcessEvent} JSON being correctly serialised by producers and
correctly parsed by monitoring topology consumers. Any mismatch would
silently corrupt state stores.

% ===========================================================================

\subsection{ProcessStateStoreIntegrationTest}
\label{it:processStateStore}

\textbf{Purpose:} Verify that the compacted-topic \texttt{ProcessStateStore}
persists and retrieves \texttt{ProcessState} records with correct
latest-wins semantics.

\textbf{Topic:} A single compacted topic (\texttt{cleanup.policy=compact}).
The store is a thin wrapper around a \texttt{KafkaProducer} (for writes) and
a \texttt{KafkaConsumer} (for reads on a compacted topic).

\textbf{Scenarios (5):}

\begin{description}

\item[shouldSaveAndLoadLatest] Saves one \texttt{ProcessState} with a known
\texttt{processInstanceId}, then loads it back. Asserts the returned state
has the same instance ID, version, variables, and token list.

\item[shouldReturnLatestAfterMultipleSaves] Saves three states for the same
key with incrementing version numbers (\texttt{ver=1,2,3}). The load
must return version~3.

\item[shouldHandleMultipleKeys] Saves three distinct keys
(\texttt{pi-a}, \texttt{pi-b}, \texttt{pi-c}) and verifies each is
independently retrievable with the correct variables.

\item[shouldReturnNullWhenKeyNotFound] Calls \texttt{loadLatest} on a
non-existent key and asserts \texttt{null} is returned within the timeout.

\item[shouldHandleTimeoutGracefully] Calls \texttt{loadLatest} with a zero
duration timeout and verifies \texttt{null} is returned cleanly (no exception).

\end{description}

\textbf{Why it matters:} The \texttt{ProcessStateStore} is used by generated
projects to snapshot process state to a compacted topic. Correct
latest-read semantics are critical for crash recovery.

% ===========================================================================

\subsection{MonitoringTopologyIntegrationTest}
\label{it:monitoringTopology}

\textbf{Purpose:} End-to-end verification of the Kafka Streams monitoring
topology against a real broker.

\textbf{Topics:} Seven topics (events, state, counts, active, latency,
trends, models) created with UUID suffixes. The topology uses
\texttt{process-events-.*} regex subscription to consume from the events
topic and materialises results into state stores queried through
\texttt{ProcessMonitoringQueryService}.

\textbf{Scenarios (5):}

\begin{description}

\item[shouldMaterializeProcessStateFromLifecycleEvents] Publishes three
events (\texttt{PROCESS\_STARTED}, \\ \texttt{ACTIVITY\_ENTERED},
\texttt{ACTIVITY\_COMPLETED}) for one instance, waits for state
materialisation, and asserts the lifecycle is \texttt{"active"} with the
correct \texttt{currentActivityId} and a populated \\
\texttt{activityDurationsMs} map.

\item[shouldProduceStateCounts] Publishes events for two instances of the
same process, then queries the count store. Asserts at least one instance
appears in the count aggregation.

\item[shouldProduceLatencySummaries] Publishes an
\texttt{ACTIVITY\_ENTERED} followed 4~minutes later by \\
\texttt{ACTIVITY\_COMPLETED} for the same activity. Queries the latency
store and asserts the returned \\ \texttt{ActivityLatencySummary} has a
\texttt{sampleCount} $\ge 1$ and \texttt{maxDurationMs} $\ge 240\,000$~ms.
This validates the entry-store-based latency computation pipeline.

\item[shouldDetectStuckInstances] Publishes an
\texttt{ACTIVITY\_ENTERED} that is never completed, then queries
\texttt{stuckInstances()} with a 1-second threshold. Asserts the instance
appears as stuck with \\ \texttt{lifecycleState} \texttt{"active"}.

\item[shouldMaterializeProcessStateFromUnkeyedLifecycleEvents] Publishes
a \texttt{PROCESS\_STARTED} event with a \texttt{null} Kafka message key.
Asserts that the monitoring topology extracts the
\texttt{processInstanceId} from the value payload and materialises state
correctly. This verifies the unkeyed-event fallback path in the topology.

\end{description}

\textbf{Why it matters:} This is the most important integration test suite.
It validates the monitoring topology's ability to materialise state, compute
latency percentiles, detect stuck instances, and handle edge cases like
unkeyed events --- all over a real Kafka broker.

% ===========================================================================

\subsection{ChaosIntegrationTest}
\label{it:chaos}

\textbf{Purpose:} Resilience verification of the monitoring topology under
crash-and-restart cycles.

\textbf{Infrastructure:} Same topic set as \texttt{MonitoringTopologyIntegrationTest}.
Uses a fixed \\ \texttt{application.id} across restarts so that Kafka Streams
rebalances and replays from the changelog topics.

\textbf{Scenarios (3):}

\begin{description}

\item[shouldRecoverStateAfterTopologyRestart] Publishes a complete task
flow (\texttt{PROCESS\_STARTED} $\to$ \\ \texttt{ACTIVITY\_ENTERED} $\to$
\texttt{ACTIVITY\_COMPLETED} $\to$ \texttt{PROCESS\_COMPLETED}), waits for
the state to reach \texttt{"completed"}, kills the topology, restarts it,
waits for re-materialisation, and asserts the recovered state matches the
pre-crash state (same \texttt{processInstanceId}, \texttt{processId},
\texttt{lifecycleState}).

\item[shouldNotDuplicateStateEntriesAfterRestart] Publishes a complete
flow, records the count-store total, kills and restarts the topology,
and asserts the count-store total is unchanged --- no spurious
duplicate events are generated by the rebalance.

\item[shouldHandleMultipleRestarts] Publishes data once, then cycles the
topology through three start-kill-restart iterations (0, 1, 2 restarts),
asserting at each iteration that the original state is still present
with the correct \texttt{lifecycleState} and \texttt{processInstanceId}.

\end{description}

\textbf{Why it matters:} Durga is designed for continuous operations;
topology restarts happen during deployments, scaling events, and failure
recovery. These tests verify that the state-store changelogs produce
correctly idempotent results.

% ===========================================================================

\subsection{E2EPipelineIntegrationTest}
\label{it:e2ePipeline}

\textbf{Purpose:} Full data pipeline end-to-end verification combining
scaffolder output validation with monitoring topology event handling for
all category-specific lifecycle events introduced in the current release.

\textbf{Fixture:} \filepath{durga-tools/src/test/resources/bpmn/e2e\_pipeline.bpmn} ---
a pipeline exercising \texttt{json-transform}, \texttt{type-coercer},
\texttt{field-router} (route), \texttt{kv-enricher} (enrich), \\
\texttt{json-schema-validator} (validate), and \texttt{window-counter}
(aggregate), with an XOR gateway branching on the \texttt{amount} field.

\textbf{Scenarios (5):}

\begin{description}

\item[shouldGeneratePluginExecutorsWithCategorySpecificEvents] Runs the
scaffolder on the E2E fixture and inspects the generated executor classes:
\begin{itemize}
\item \texttt{RouteDecisionPluginExecutor.java} must contain
\texttt{ProcessEvent.EventType.\\GATEWAY\_TAKEN} (route category).
\item \texttt{ValidateHighValuePluginExecutor.java} must contain
\texttt{ProcessEvent.EventType.\\ACTIVITY\_ESCALATED},
\texttt{ProcessEvent.Status.ESCALATED}, and error code \\
\texttt{"VALIDATION\_FAILED"} (validate category).
\item \texttt{AggregateHighValuePluginExecutor.java} must contain
\texttt{ProcessEvent.EventType.\\ACTIVITY\_ENTERED} (aggregate category,
null-output path).
\end{itemize}

\item[shouldHandleGatewayTakenEvents] Publishes
\texttt{PROCESS\_STARTED} $\to$ \texttt{ACTIVITY\_ENTERED} $\to$ \\
\texttt{GATEWAY\_TAKEN} for \texttt{route\_decision} and verifies the
monitoring state view records a duration for the routed activity with
lifecycle \texttt{"active"}.

\item[shouldHandleActivityEscalatedEvents] Publishes
\texttt{PROCESS\_STARTED} $\to$ \texttt{ACTIVITY\_ENTERED} $\to$
\texttt{ACTIVITY\_ESCALATED} for \texttt{validate\_high\_value} and
verifies the monitoring topology records the escalated activity's
duration correctly.

\item[shouldHandleActivityEnteredThenCompletedEvents] Publishes
\texttt{PROCESS\_STARTED} $\to$ \\ \texttt{ACTIVITY\_ENTERED} $\to$
\texttt{ACTIVITY\_COMPLETED} for \texttt{aggregate\_high\_value} (30~s
gap) and asserts that the latency query returns an
\texttt{ActivityLatencySummary} with $\ge 1$ sample and \\
\texttt{maxDurationMs} $\ge 30\,000$~ms. This validates the entry-store
latency pipeline for aggregate window-start events.

\item[shouldHandleMultipleEventTypesInLifecycle] Simulates a
full pipeline trace across five activities with mixed event types:
\begin{itemize}
\item \texttt{transform\_order}: entered $\to$ completed
\item \texttt{coerce\_types}: entered $\to$ completed
\item \texttt{route\_decision}: entered $\to$ \texttt{GATEWAY\_TAKEN}
\item \texttt{enrich\_high\_value}: entered $\to$ completed
\item \texttt{validate\_high\_value}: entered $\to$
\texttt{ACTIVITY\_ESCALATED}
\end{itemize}
Asserts all five activities appear in the state view's duration map and
the lifecycle state is \texttt{"active"}. Each test method uses an
isolated \texttt{processId} to prevent cross-test contamination in the
shared latency store.

\end{description}

\textbf{Why it matters:} This test closes the gap between the scaffolder's
code-generation contract and the monitoring topology's runtime behaviour.
It proves that the category-specific event types (\texttt{GATEWAY\_TAKEN},
\texttt{ACTIVITY\_ESCALATED}, \texttt{ACTIVITY\_ENTERED}) are correctly
generated by the scaffolder AND correctly consumed and materialised by
the monitoring engine --- end to end, against a live Kafka broker.

% ===========================================================================

\subsection{FaultDetectionTest (unit)}
\label{it:faultDetection}

\textbf{Purpose:} Unit-test the fault detection algorithm across all three
syndromes without requiring a Kafka broker. The test uses in-memory
\texttt{MapKVStore} adapters that implement the \texttt{KeyValueStore}
interface.

\textbf{Class:} \texttt{FaultDetectionTest} (14 scenarios, no Docker required).

\textbf{Scenarios (14):}

\begin{description}

\item[shouldFireHardErrorOnFirstOccurrence] Config with scope
(\texttt{processId="order-proc"}, \\ \texttt{activityId="validate\_order"}),
\texttt{HARD\_ERROR} syndrome, \texttt{CRITICAL} severity. Asserts alarm
fires with correct message rendering, scope fields, and count=1.

\item[shouldFireHardErrorRepeatedly] Wildcard-scoped config (both
\texttt{processId} and \texttt{activityId} null) fires on every matching
\texttt{PROCESS\_FAILED} event regardless of process or activity.

\item[shouldNotFireHardErrorForNonMatchingEventType] Config targeting
\texttt{ACTIVITY\_ESCALATED} \\ ignores a \texttt{PROCESS\_FAILED} event.

\item[shouldMatchByProcessOnly] Config scoped to \texttt{"order-proc"}
fires for events in that process, ignores events in \texttt{"other-proc"}.

\item[shouldMatchByActivityOnly] Config scoped to
\texttt{"validate\_order"} fires for that activity in any process, ignores
other activities.

\item[shouldCountAndFireWhenThresholdExceeded] \texttt{COUNTED} syndrome
with threshold=3. First three events increment count silently; fourth
crosses threshold and fires alarm; fifth is suppressed (single-fire).

\item[shouldReArmAfterCountDropsBelowThreshold] After an alarm fires,
manually resets count below threshold and clears the alarm marker. The
next event that crosses threshold fires a new alarm.

\item[shouldFireWhenWindowCountExceedsThreshold] \texttt{SLIDING\_WINDOW}
with threshold=2 and 60-second window. First three events accumulate;
third crosses threshold and fires; fourth suppressed (single-fire).

\item[shouldPruneOldTimestampsInWindow] With a 100-ms window, inserts
four events, waits 150 ms, then adds a fifth. The four old timestamps
are pruned; only the one new event remains in the window, so no alarm
fires.

\item[shouldHandleMultipleConfigs] A \texttt{HARD\_ERROR} and a
\texttt{COUNTED} config both match the same event. The hard-error
fires immediately; the counted config correctly matches the event's
scope and type.

\item[shouldHandleEscalatedEvents] Config targeting
\texttt{ACTIVITY\_ESCALATED} fires on an escalated event with correct
syndrome and config ID.

\item[shouldBuildTopologyWithRegexEventsAndModelRegistry] Verifies that the
fault detection topology subscribes to lifecycle event topics by regex and
includes the process-model registry input needed for BPMN alarm configuration.

\item[shouldRejectSlidingWindowWithoutDuration] Constructor validation
rejects a \texttt{SLIDING\_WINDOW} config with \texttt{null}
\texttt{windowDuration}.

\item[shouldRejectCountedWithoutThreshold] Constructor validation
rejects a \texttt{COUNTED} config with \\ \texttt{threshold=0}.

\end{description}

\textbf{Why it matters:} The fault detection algorithm is the core
intelligence that distinguishes between hard errors, counted thresholds,
and time-windowed bursts. Getting the de-duplication and re-arm logic
wrong would cause either alarm floods or silent failures.

% ===========================================================================

\subsection{AlarmStateViewTest and AlarmStateTopologyTest (unit)}
\label{it:alarmState}

\textbf{Purpose:} Unit-test the alarm-state read model that folds emitted
\texttt{AlarmEvent} records into queryable \texttt{AlarmStateView} records for
the dashboard and HTTP API.

\textbf{Classes:} \texttt{AlarmStateViewTest} (4 scenarios) and
\texttt{AlarmStateTopologyTest} (1 scenario), no Docker required.

\textbf{Scenarios (5):}

\begin{description}

\item[shouldBuildStateFromAlarmEvent] Converts a single \texttt{AlarmEvent} into
an \texttt{ACTIVE} state view with the expected key, timestamps, fire count, and
message.

\item[shouldMergeRepeatedAlarmEventsAndKeepLatestDetails] Folds repeated firings
for the same \\ process/activity/config into one view, increments \texttt{fireCount},
keeps the strongest severity, and records the latest firing details.

\item[shouldNotLetLateOlderEventOverwriteLatestDetails] Ensures an older event
that arrives late does not overwrite the newest alarm ID, instance ID, message,
or count.

\item[shouldUseWildcardActivityForProcessWideAlarms] Verifies that process-wide
alarms use \texttt{*} as the activity component in the alarm-state key.

\item[shouldBuildAlarmStateStoreTopology] Verifies that the alarm-state topology
consumes \texttt{fault-alarms} and materializes \texttt{alarm-state-store}.

\end{description}

\textbf{Why it matters:} Fault detection emits a stream of alarm firings, but
operators need a stable read model. These tests protect the fold-up semantics
used by \texttt{/api/alarms}, per-process alarm queries, instance alarm details,
and dashboard overlays.

% ===========================================================================

\subsection{Infrastructure}

The Docker-backed suites share a common base class:

\textbf{KafkaIntegrationTestBase} (\filepath{src/test/java/org/gautelis/}
\filepath{durga/KafkaIntegrationTestBase.java}) uses the Testcontainers
library to start a single Confluent CP Kafka 7.8.0 broker in KRaft mode
(no ZooKeeper). The broker listens on a random free host port and
advertises \fbox{\texttt{localhost:<port>}}. Docker availability is
auto-detected; if Docker is missing, tests are skipped unless \\
\texttt{DURGA\_REQUIRE\_DOCKER\_TESTS=true} is set.

\begin{itemize}
\item \textbf{Environment variables:}
\texttt{DURGA\_DOCKER\_HOST\_IP} (override advertised listener host), \\
\texttt{DURGA\_REQUIRE\_DOCKER\_TESTS} (fail if Docker unavailable),
\texttt{DURGA\_IT\_TIMEOUT\_SECONDS} (override 60~s default).
\item \textbf{System properties:}
\texttt{durga.it.timeout.seconds} (same as env var, takes precedence), \\
\texttt{DOCKER\_HOST} (Docker socket URI, auto-detected if unset).
\end{itemize}

Each integration test class creates topics with random UUID suffixes in
its \texttt{@BeforeClass} and tears down the broker in \texttt{@AfterClass}.
Test methods create their own \texttt{KafkaProducer}, \texttt{KafkaConsumer},
or \texttt{KafkaStreams} instances in \texttt{@Before} and close them in
\texttt{@After}.

\begin{itemize}
\item \textbf{Run all integration tests:}
\texttt{mvn test -Dtest='*IntegrationTest'}

\item \textbf{Run a specific suite:}
\texttt{mvn test -Dtest='E2EPipelineIntegrationTest'}

\item \textbf{Require Docker (fail if unavailable):}
\texttt{mvn test -Ddurga.integration.requireDocker=true \\-Dtest='*IntegrationTest' }
\end{itemize}
